\chapter{Source Code of the Fergana Transport Bot}
\label{app:source-code}

This appendix contains the main source code files of the Fergana Transport
Bot. The code is presented in two parts: the database initialization
script (\texttt{init\_db.py}) and a representative excerpt of the bot
logic (\texttt{bot.py}). Repeated handler code has been shortened for
readability. The full and most recent code is available in the project
repository on GitHub.

\section{Database Initialization Script (\texttt{init\_db.py})}
\label{app:init-db}

This script creates the SQLite tables defined in
Listing~\ref{lst:schema}, clears any existing data, imports the route
data from the \texttt{data/} folder (HTML pages saved from Yandex Maps),
and inserts demo taxi services. The script is described in
Section~\ref{subsec:data}.

\begin{lstlisting}[language=Python, caption={Full source of
    \texttt{init\_db.py}}, label={lst:init-db}]
import os
import re
import html
import sqlite3


DB_NAME = "transport.db"
DATA_DIR = "data"


def get_connection():
    return sqlite3.connect(DB_NAME)


def create_tables(conn):
    cursor = conn.cursor()

    cursor.execute("""
        CREATE TABLE IF NOT EXISTS routes (
            id INTEGER PRIMARY KEY AUTOINCREMENT,
            number TEXT NOT NULL UNIQUE,
            name TEXT,
            type TEXT,
            description TEXT
        )
    """)

    cursor.execute("""
        CREATE TABLE IF NOT EXISTS stops (
            id INTEGER PRIMARY KEY AUTOINCREMENT,
            name TEXT NOT NULL UNIQUE,
            latitude REAL,
            longitude REAL,
            district TEXT
        )
    """)

    cursor.execute("""
        CREATE TABLE IF NOT EXISTS route_stops (
            id INTEGER PRIMARY KEY AUTOINCREMENT,
            route_id INTEGER NOT NULL,
            stop_id INTEGER NOT NULL,
            stop_order INTEGER NOT NULL,
            FOREIGN KEY (route_id) REFERENCES routes(id),
            FOREIGN KEY (stop_id) REFERENCES stops(id)
        )
    """)

    cursor.execute("""
        CREATE TABLE IF NOT EXISTS taxi_services (
            id INTEGER PRIMARY KEY AUTOINCREMENT,
            name TEXT NOT NULL,
            phone TEXT,
            working_hours TEXT
        )
    """)

    cursor.execute("""
        CREATE TABLE IF NOT EXISTS users (
            id INTEGER PRIMARY KEY AUTOINCREMENT,
            telegram_id INTEGER NOT NULL UNIQUE,
            username TEXT,
            registered_at TEXT DEFAULT CURRENT_TIMESTAMP
        )
    """)

    conn.commit()


def clear_data(conn):
    cursor = conn.cursor()
    cursor.execute("DELETE FROM route_stops")
    cursor.execute("DELETE FROM routes")
    cursor.execute("DELETE FROM stops")
    cursor.execute("DELETE FROM taxi_services")
    conn.commit()


def clean_text(value):
    value = html.unescape(value)
    value = re.sub(r"\s+", " ", value)
    return value.strip()


def extract_route_number(filename, content):
    match = re.search(
        r'masstransit-card-header-view__number[^>]*title="([^"]+)"',
        content)
    if match:
        return clean_text(match.group(1))
    match = re.search(r"bus(\d+)", filename.lower())
    if match:
        return match.group(1)
    return None


def extract_essential_stops(content):
    pattern = r'masstransit-card-header-view__essential-stop[^>]*>(.*?)</'
    found = re.findall(pattern, content, flags=re.DOTALL)
    stops = []
    for item in found:
        item = re.sub(r"<.*?>", "", item)
        item = clean_text(item)
        if item:
            stops.append(item)
    return stops


def extract_all_stops(content):
    stops = []
    aria_matches = re.findall(
        r'<a[^>]*class="[^"]*masstransit-legend-group-view__item-link'
        r'[^"]*"[^>]*aria-label="([^"]+)"',
        content, flags=re.DOTALL
    )
    for item in aria_matches:
        item = clean_text(item)
        if item and item not in stops:
            stops.append(item)
    return stops


def insert_route(conn, number, essential_stops):
    cursor = conn.cursor()
    if len(essential_stops) >= 2:
        name = f"{essential_stops[0]} - {essential_stops[-1]}"
        description = (
            f"Route {number}: {essential_stops[0]} "
            f"to {essential_stops[-1]}"
        )
    else:
        name = f"Route {number}"
        description = f"Bus route {number} in Fergana"
    cursor.execute("""
        INSERT OR IGNORE INTO routes (number, name, type, description)
        VALUES (?, ?, ?, ?)
    """, (number, name, "bus", description))
    conn.commit()
    cursor.execute("SELECT id FROM routes WHERE number = ?", (number,))
    return cursor.fetchone()[0]


def insert_stop(conn, stop_name):
    cursor = conn.cursor()
    cursor.execute("""
        INSERT OR IGNORE INTO stops (name, latitude, longitude, district)
        VALUES (?, NULL, NULL, NULL)
    """, (stop_name,))
    conn.commit()
    cursor.execute("SELECT id FROM stops WHERE name = ?", (stop_name,))
    return cursor.fetchone()[0]


def insert_route_stop(conn, route_id, stop_id, stop_order):
    cursor = conn.cursor()
    cursor.execute("""
        INSERT INTO route_stops (route_id, stop_id, stop_order)
        VALUES (?, ?, ?)
    """, (route_id, stop_id, stop_order))
    conn.commit()


def add_demo_taxi_services(conn):
    cursor = conn.cursor()
    taxi_data = [
        ("Taxi information", "Not available yet", "24/7"),
        ("Administrator can add taxi services later",
         "Not available yet", "24/7")
    ]
    for name, phone, hours in taxi_data:
        cursor.execute("""
            INSERT INTO taxi_services (name, phone, working_hours)
            VALUES (?, ?, ?)
        """, (name, phone, hours))
    conn.commit()


def import_routes_from_files(conn):
    if not os.path.exists(DATA_DIR):
        print("Folder 'data' was not found.")
        return
    files = sorted([
        file for file in os.listdir(DATA_DIR)
        if file.lower().startswith("bus")
        and file.lower().endswith(".txt")
    ])
    for filename in files:
        path = os.path.join(DATA_DIR, filename)
        with open(path, "r", encoding="utf-8") as file:
            content = file.read()
        route_number = extract_route_number(filename, content)
        essential_stops = extract_essential_stops(content)
        all_stops = extract_all_stops(content)
        if not route_number or not all_stops:
            continue
        route_id = insert_route(conn, route_number, essential_stops)
        for index, stop_name in enumerate(all_stops, start=1):
            stop_id = insert_stop(conn, stop_name)
            insert_route_stop(conn, route_id, stop_id, index)
        print(f"Imported route {route_number}: "
              f"{len(all_stops)} stops")


def main():
    conn = get_connection()
    create_tables(conn)
    clear_data(conn)
    import_routes_from_files(conn)
    add_demo_taxi_services(conn)
    conn.close()
    print("Database created successfully.")


if __name__ == "__main__":
    main()
\end{lstlisting}

\section{Bot Application (\texttt{bot.py}, excerpt)}
\label{app:bot-py}

This listing shows the main structure of the bot application. The handlers
for \texttt{/start}, \texttt{/help}, \texttt{/marshrut}, \texttt{/poisk},
\texttt{/taxi}, and \texttt{/dobavit} follow the same pattern; only the
two most important ones are shown in full to save space.

\begin{lstlisting}[language=Python, caption={Excerpt from
    \texttt{bot.py}}, label={lst:bot-py}]
import asyncio
import os
import sqlite3
from dotenv import load_dotenv
from aiogram import Bot, Dispatcher, types
from aiogram.filters import Command, CommandStart

load_dotenv()
BOT_TOKEN = os.getenv("BOT_TOKEN")
DB_NAME = "transport.db"

bot = Bot(token=BOT_TOKEN)
dp = Dispatcher()


def get_stops_by_route(route_number):
    conn = sqlite3.connect(DB_NAME)
    cursor = conn.cursor()
    cursor.execute("""
        SELECT s.name FROM stops s
        JOIN route_stops rs ON rs.stop_id = s.id
        JOIN routes r ON r.id = rs.route_id
        WHERE r.number = ?
        ORDER BY rs.stop_order
    """, (route_number,))
    rows = cursor.fetchall()
    conn.close()
    return [row[0] for row in rows]


def find_route_between(stop_a, stop_b):
    conn = sqlite3.connect(DB_NAME)
    cursor = conn.cursor()
    cursor.execute("""
        SELECT DISTINCT r.number, r.name FROM routes r
        JOIN route_stops rs1 ON rs1.route_id = r.id
        JOIN stops s1 ON s1.id = rs1.stop_id
        JOIN route_stops rs2 ON rs2.route_id = r.id
        JOIN stops s2 ON s2.id = rs2.stop_id
        WHERE s1.name LIKE ? AND s2.name LIKE ?
          AND rs1.stop_order < rs2.stop_order
        LIMIT 5
    """, (f"%{stop_a}%", f"%{stop_b}%"))
    rows = cursor.fetchall()
    conn.close()
    return rows


@dp.message(CommandStart())
async def handle_start(message: types.Message):
    text = (
        "\U0001F44B Welcome to Fergana Transport Bot!\n\n"
        "This bot helps you find public transport information "
        "in Fergana.\n\n"
        "Available commands:\n"
        "/marshrut 1 - show route stops\n"
        "/marshrut 2 - show route stops\n"
        "/marshrut 4 - show route stops\n"
        "/poisk Stop A - Stop B - find route between stops\n"
        "/taxi - show taxi information\n"
        "/help - show help"
    )
    await message.answer(text)


@dp.message(Command("poisk"))
async def handle_poisk(message: types.Message):
    args = message.text.split(maxsplit=1)
    if len(args) < 2 or "-" not in args[1]:
        await message.answer(
            "Please use this format:\n"
            "/poisk Stop A - Stop B\n\n"
            "Example:\n"
            "/poisk Istiklol - Aeroport"
        )
        return
    parts = args[1].split("-", maxsplit=1)
    stop_a = parts[0].strip()
    stop_b = parts[1].strip()
    results = find_route_between(stop_a, stop_b)
    if not results:
        await message.answer("No route found between these stops.")
        return
    text = (
        f"Search result:\n\n"
        f"From: {stop_a}\n"
        f"To: {stop_b}\n\n"
    )
    for number, name in results:
        text += f"\U0001F68C Route {number}: {name}\n"
    await message.answer(text)


# Other handlers (/help, /marshrut, /taxi, /dobavit)
# follow the same pattern.

async def main():
    await dp.start_polling(bot)


if __name__ == "__main__":
    asyncio.run(main())
\end{lstlisting}